\subsection{Proof of Theorem~\ref{thm:14.6i}}\label{sec:16}
We first introduce some notation in preparation for the proofs to come.
\begin{notation}\label{not:rho}
	\begin{enumerate}[label=\textup{(\roman*)}]
		\item \label{561} Let $\lambda$ be a partition, $n\in\N_0$ and $\phi$ be a virtual character of $S_n$.
		We define
		\[ \phi/\chi^\lambda\coloneqq  \sum_{\mu\vdash n}\langle \phi,\chi^\mu\rangle\cdot\chi^{\mu/\lambda} \]
		where $\chi^{\mu/\lambda}=0$ if $\lambda\not\subseteq\mu$.
		
		\item \label{562} For $m,n\in\N$ and $\alpha/\beta$ a skew shape of size $n$, define \[
		\rho^{\alpha/\beta}_m \coloneqq  \cX(\triv_{S_m};\chi^{\alpha/\beta})\up_{S_m\wr S_n}^{S_{mn}}.\] 
		
		\item \label{563} For $\phi\in\Ch(S_{n_1})$ and $\theta\in\Ch(S_{n_2})$, define
		\[ \phi\boxtimes\theta\coloneqq  (\phi\times\theta)\up_{S_{n_1}\times S_{n_2}}^{S_{n_1+n_2}}. \]
		
		\item \label{564} Let $S_\lambda$ denote the Young subgroup $S_{\lambda_1}\times\cdots\times S_{\lambda_{l(\lambda)}}$ of $S_n$ and let
		\[ \zeta^\lambda\coloneqq \triv_{S_\lambda}\up^{S_n} = \chi^{(\lambda_1)}\boxtimes\cdots\boxtimes \chi^{(\lambda_{l(\lambda)})} \]
		denote the character of the permutation module $M^\lambda$. (Note $\zeta^\mu=\zeta^\lambda$ if $\mu$ is a composition of $n$ with the same parts as $\lambda$ but in a different order.) 
	\end{enumerate}
\end{notation}

Recall that the irreducible decomposition of such a permutation character $\zeta^\lambda$ is described by Young's Rule \cite[2.8.5]{JK}. Equivalently,
\begin{equation}\label{eqn:kostka}
	\zeta^\lambda = \sum_{\gamma\vdash n} K_{\gamma,\lambda}\cdot \chi^\gamma
\end{equation}
where the Kostka number $K_{\gamma,\lambda} = \langle \zeta^\lambda,\chi^\gamma\rangle=c^\gamma_{(\lambda_1),\dotsc,(\lambda_{l(\lambda)})}$ equals the number of semistandard Young tableaux of shape $\gamma$ and content $\lambda$.
In particular, we therefore have
\begin{equation}\label{eqn:15.2}
	\sum_{\gamma\vdash n} K_{\gamma,\lambda}\cdot \rho^\gamma_m = \cX(\triv_{S_m}; \zeta^\lambda)\up_{S_m\wr S_n}^{S_{mn}} = \rho^{(\lambda_1)}_m\boxtimes\cdots\boxtimes\rho^{(\lambda_{l(\lambda)})}_m.
\end{equation}

We now precisely identify the $i=k$ term of Theorem~\ref{thm:14.6i}, in Theorem~\ref{thm:14.6ii}. We then deduce Theorem~\ref{thm:14.6i} from Theorem~\ref{thm:14.6ii}. Following that, we prove Theorem~\ref{thm:14.6ii}, during the course of which we will also prove Theorem B, which has been numbered as Theorem~\ref{thm:16.7} in this section for ease of reference.
